\documentclass[11pt,letterpaper]{article}
\usepackage[T1]{fontenc}
\usepackage{lmodern,amsmath,amssymb,amsthm,microtype}
\usepackage[margin=1in]{geometry}
\usepackage[hidelinks]{hyperref}
\newtheorem{theorem}{Theorem}
\title{Four-point extensions of \(G_{449}\) are five-colorable}
\author{Aafi Khan Malek\\ \small Independent researcher, Ahmedabad, India\\
\small\href{mailto:aafimalek2032@gmail.com}{aafimalek2032@gmail.com}}
\date{28 November 2026}
\begin{document}
\maketitle
Let \(P\) be the 449 points from Heule and Przybocki's
\href{https://arxiv.org/abs/2611.28518}{\(G_{449}\) construction}.
We give an exact finite verification for four arbitrary added points.
A six-chromatic induced extension of this fixed point set would
need at least five new points, and if it has exactly five, all five
must lie in the denominator-seven cyclotomic grid. These statements
do not change the bounds for coloring the whole plane.

\begin{theorem}
For any set \(X\) of at most four points of the Euclidean plane,
the induced unit-distance graph on \(P\cup X\) is five-colorable.
\end{theorem}

\section{Reduction}
Use the grid \(L=\tfrac17\mathbb Z[\zeta]\) for
\(\zeta=e^{2\pi i/42}\). The preceding
\href{https://github.com/Aafimalek/novel-discoveries/blob/6caa1be3f702e8f1b1b56ef5497eea6e885ccfd7/hadwiger-nelson/verification/g449/prime-seven-catalogue/paper/three-point-extensions.pdf}
{three-point verification} gives five-colorability after deleting any
one of four distinct new points, and proves that an off-grid point
has at most two neighbors in \(P\). In a failing extension every new
vertex must have degree at least five, since we could otherwise color
it greedily after coloring the other three.

\section{An off-grid point}
Fix an off-grid new vertex \(x\). It must have exactly two base
neighbors and be adjacent to the three other new vertices. There
can be at most two off-grid new vertices: three would force four
pairwise equidistant planar points. Two grid neighbors of \(x\)
cannot be adjacent. Indeed the two points at unit distance from
both endpoints of a grid unit edge belong to the grid, because
rotating a grid direction by \(\pm60^\circ\) multiplies it by
\(\zeta^{\pm7}\). In particular the two base neighbors of \(x\)
are nonadjacent. At least one other new point \(y\) is on the grid;
it has at most two new neighbors because its grid companions are
also adjacent to \(x\), so it has at least three base neighbors.

The mixed-triple calculation rules out this configuration.
Consider all 238,960,162 triples consisting of two nonadjacent
base vertices and one distinct grid neighbor of at least three
base vertices. The latter points are exactly 2,429 outside grid
points listed by the complete catalogue of unit directions.
After clearing the common denominator seven, put
\(u=B-A\), \(v=C-A\), \(U=u\bar u\), \(V=v\bar v\),
\(W=(u-v)\overline{(u-v)}\) and \(D=u\bar v-\bar u v\)
for such a triple. Unit circumradius implies
\(UVW+49D^2=0\). This necessary condition may be tested
at reciprocal primitive 42nd roots modulo each of
\(43,127,211,337\). Exactly 567,847 triples survive, and
the checker exhibits a grid center for each from the unit
directions. A triple of distinct points on a circle has a
unique circumcenter. Thus none has an off-grid center, as
required.

\section{Four grid points}
A new vertex has at most three new neighbors, so only grid points
with at least two base neighbors are candidates. Translating
the complete unit-direction set from the base enumerates exactly
7,056 such candidates. Induced adjacency on these points is
again an exact unit-direction lookup.

The bank contains 36 proper colorings of \(P\): the 35 from
the three-point calculation and the restriction of its
503-vertex witness. For each base coloring, the available list of
a candidate consists of its colors absent among base neighbors.
A component of the new graph can be assigned colors exactly when
these lists admit a proper coloring of that component. The checker
uses the exact list test on each component, combines the
36-bit success masks by intersection, and verifies all instances
with empty intersection by explicit color witnesses.

Only new sets satisfying the degree restriction need be checked:
a vertex in a component of size \(s\) needs as many base neighbors
as are required to bring its actual total degree to five.
Enumerate connected components by extending a smaller connected
set with an adjacent vertex, deduplicate, and test their internal
degrees. For disconnected sets, combine the component masks by
size partition and reject combinations with repeated vertices or
edges across components. This includes all cases of the
partitions \(1+1+1+1\), \(2+1+1\), \(2+2\), \(3+1\), and \(4\).

In the candidate graph there are 33,908 edges and 115,395
connected four-element sets that pass the degree restriction.
The bank leaves 8,629 distinct four-sets. In the component
order above, the counts are 6,218, 2,322, 22, 66, and 1.
The union of these sets consists of just 286 new vertices.
With (P) they form a 735-vertex induced graph with 4,107
edges. A single stored five-coloring of this graph, verified
on each edge, covers every remaining four-set. An independent Laurent-polynomial
calculation tests all 269,745 pairwise distances in this graph
and verifies the coloring. This proves the theorem.

\section{Five points with at least one off-grid point}
\begin{theorem}
If at least one of five distinct new points is outside \(L\), the
induced unit-distance graph on their union with \(P\) is five-colorable.
\end{theorem}
 Any two points have at most two common unit
neighbors; two adjacent points of $L$ have both common unit neighbors
in $L$ (rotate by $\pm60^\circ$). Hence the grid neighbors of any
off-grid point are independent. We use two exact certificates:
an off-grid point with two base neighbors has no grid neighbor of
base degree at least three, and an off-grid point with one base
neighbor cannot have two grid neighbors of base degree at least four.
For the latter, test one base point and two of the 1,043 grid points
of base degree at least four; exclude triples containing a unit
edge. Of the remaining 237,808,116 triples, precisely 353,150 pass
the necessary radius congruences at 43, 127, 211, and 337. Each
survivor has a grid circumcenter by exact direction lookup, so
uniqueness excludes an off-grid center.

In a failing five-point extension each new vertex has total degree
at least five, by deletion and greedy extension. An off-grid vertex
has at most two base neighbors, so it has at least three new
neighbors; if it has exactly three, it cannot meet any new grid
vertex of base degree at least three. Consider the number of
off-grid vertices.
\begin{itemize}
\item \emph{One.} If it has three new neighbors, each has base degree
at most two and can see only it and the fourth grid vertex among
the new points. If it has four, they are independent and all have
base degree at least four, contradicting the second certificate
because it must also have a base neighbor.
\item \emph{Two.} If nonadjacent, both meet all three grid vertices.
Those are independent with base degree at least three, a contradiction
at either off-grid point. If adjacent and one meets all three, the
grid vertices are again independent with base degree at least three,
so the other must also meet all three; they cannot have three
common neighbors. Otherwise each meets exactly two grid vertices,
which must have base degree at most two. If the two neighbor sets
differ, a grid vertex in just one set has at most two new neighbors
(one off-grid point and at most one grid point). If the sets
coincide, the two common grid neighbors each need one more edge
because their base degrees are at most two. The third grid vertex
would therefore be a common neighbor, but both common unit
neighbors are already off-grid.
\item \emph{Three.} The off-grid graph is a path or a triangle. In
the path case its leaves meet both grid vertices, forcing their
base degrees to be at most two. These grid vertices cannot be adjacent,
so they must both meet all three off-grid vertices, impossible.
For a triangle, the two grid vertices must between them meet
all three off-grid vertices, and each meets at most two.
If both meet two they are distinct exterior equilateral apices,
of forms $a+b-c$ and $b+c-a$ up to relabeling; they are distance two
apart, so both have base degree at least three. Some off-grid
vertex meets just one grid vertex, a contradiction there.
Otherwise the two grid vertices meet disjoint sets of sizes one and
two. The grid vertex meeting just one off-grid vertex has
base degree at least three, again contradicting the restriction
at its sole off-grid neighbor.
\item \emph{Four or five.} With four, their graph has minimum degree
two, hence is a cycle or a diamond (a unit $K_4$ is impossible).
Two nonadjacent degree-two vertices already share two neighbors
and would both need adjacency to the grid vertex. With five, the
graph has minimum degree three; a missing edge would have three
common neighbors, and no missing edge gives a $K_4$.
\end{itemize}
Thus every five-point extension containing an off-grid point
is five-colorable. The case of five grid points is not checked here.

\paragraph{Reproducibility and scope.}
The companion scripts use pinned base data, exact integer
coordinates, and explicit color vectors. Neither a SAT solve
nor a floating-point test is needed in the final replay.
Neither extension theorem is claimed to be checked in Lean.

\paragraph{Contributions.}
Aafi Khan Malek selected and directed the research. ChatGPT assisted
with proof development, software and writing. Codex was not used.
This research note has not been peer reviewed.
\end{document}
